\section{Introduction}
Employee attrition describes employees leaving an organisation. In this application, recorded employee characteristics are used to distinguish employees who leave from those who remain, giving a supervised binary classification task. The analysis concerns associations in the available data rather than explanations of individual decisions to leave.

Feature selection can reduce dimensionality, remove irrelevant or redundant information and make the retained employee characteristics easier to interpret. These potential benefits need to be assessed against predictive performance: a smaller representation need not preserve all useful information.

This integrated study evaluates five PA1 classifiers and four PA2 probabilistic classifiers under four feature-selection strategies: all valid predictors, univariate filtering, multivariate filtering and wrapper selection. PA2 also evaluates three metaclassifiers. The research question is: \emph{Which employee characteristics repeatedly support attrition prediction, and how do classifier choice and feature selection affect predictive performance, minority-class sensitivity and probability quality?} Dimensionality and model complexity inform the comparison.
